\section{Listas, tuplas e dicionários}

\begin{frame}[fragile,t]{8. Listas}
O que é? Listas armazenam vários elementos em uma única estrutura. São mutáveis, isto é, podem ser alteradas depois de criadas, e são representadas por colchetes.\par
\begin{visibleenv}<2->
\EMASCcodelabel{Estrutura do comando}
\begin{lstlisting}
nome_da_lista = [elemento1, elemento2, elemento3]
\end{lstlisting}
\end{visibleenv}
\end{frame}

\begin{frame}[fragile,t]{8. Listas / Exemplo}
\EMASCcodelabel{Exemplo prático}
\begin{lstlisting}
frutas = ["Maçã", "Banana", "Uva"]

print(frutas)
print(frutas[0])   # Maçã
print(frutas[1])   # Banana
\end{lstlisting}
\begin{visibleenv}<2->
\begin{EMASCnote}
Entenda o exemplo: Os índices começam em zero: o índice 0 acessa o primeiro elemento, o índice 1 acessa o segundo e assim por diante.
\end{EMASCnote}
\end{visibleenv}
\end{frame}

\begin{frame}[fragile,t]{8. Listas / Operações comuns}
\EMASCcodelabel{Exemplo prático}
\begin{lstlisting}
frutas.append("Laranja")    # Adiciona ao final
frutas.remove("Banana")     # Remove pelo valor
frutas.pop()                # Remove o último elemento
print(len(frutas))          # Quantidade de elementos
\end{lstlisting}
\end{frame}

\begin{frame}[fragile,t]{9. Tuplas}
O que é? Tuplas armazenam uma sequência de elementos e, ao contrário das listas, são imutáveis: não permitem alterar, adicionar ou remover elementos diretamente após a criação. São normalmente representadas por parênteses.\par
\begin{visibleenv}<2->
\EMASCcodelabel{Estrutura do comando}
\begin{lstlisting}
nome_da_tupla = (elemento1, elemento2, elemento3)
\end{lstlisting}
\end{visibleenv}
\end{frame}

\begin{frame}[fragile,t]{9. Tuplas / Exemplo}
\EMASCcodelabel{Exemplo prático}
\begin{lstlisting}
dias = ("Segunda", "Terça", "Quarta", "Quinta")

print(dias)
print(dias[0])
print(dias[2])
\end{lstlisting}
\begin{visibleenv}<2->
\begin{EMASCnote}
Entenda o exemplo: Assim como nas listas, os índices começam em zero. Uma tupla de um único elemento precisa de vírgula: \texttt{(5,)}.
\end{EMASCnote}
\end{visibleenv}
\end{frame}

\begin{frame}[fragile,t]{10. Dicionários}
O que é? Dicionários armazenam informações em pares de chave e valor. Cada chave identifica o valor associado. São representados por chaves.\par
\begin{visibleenv}<2->
\EMASCcodelabel{Estrutura do comando}
\begin{lstlisting}
nome_do_dicionario = {
    "chave1": valor1,
    "chave2": valor2
}
\end{lstlisting}
\end{visibleenv}
\end{frame}

\begin{frame}[fragile,t]{10. Dicionários / Exemplo}
\EMASCcodelabel{Exemplo prático}
\begin{lstlisting}
aluno = {
    "nome": "Carlos",
    "idade": 20,
    "curso": "Informática"
}

print(aluno["nome"])
print(aluno["idade"])
print(aluno["curso"])
\end{lstlisting}
\begin{visibleenv}<2->
\begin{EMASCnote}
Entenda o exemplo: As chaves são nome, idade e curso. Os valores associados são Carlos, 20 e Informática.
\end{EMASCnote}
\end{visibleenv}
\end{frame}

\begin{frame}[fragile,t]{10. Dicionários / Operações comuns}
\EMASCcodelabel{Exemplo prático}
\begin{lstlisting}
aluno["nota"] = 9.0         # Adiciona uma chave
aluno["idade"] = 21         # Atualiza um valor
aluno.pop("curso")          # Remove a chave e seu valor
print(aluno)
\end{lstlisting}
\end{frame}
